% GENERATED FILE. Edit canonical lessons, then run npm run generate:books.
% canonical-source-hash: fnv1a64:396c0ba212fdca28
% canonical-lessons: JA-C202-jidousha, JA-C202-unten, JA-C202-untenshu, JA-C202-tomeru, JA-C202-chuushajou

\chapter{Cars, Drivers and Car Parks}
\label{ch:ja-cars-drivers-and-car-parks}
\hlchaptermodality{202}

\begin{chapteropening}
\textbf{By the end of this chapter:} \emph{I can talk about cars, driving and parking.}

Complete the last lesson of chapter 202: I can talk about cars, driving and parking.
\end{chapteropening}

\section[\texorpdfstring{\ja{じどうしゃ} (jidōsha)}{じどうしゃ (jidōsha)}]{\ja{じどうしゃ} (jidōsha) — an automobile; a motor car}

\label{lesson:JA-C202-jidousha}



\begin{quote}
Before the new one: say the Japanese for to rush out, then the Japanese for to bump into.
\end{quote}

\subsection*{You'll want to know: \ja{じどうしゃ}}
\textbf{\ja{じどうしゃ}} — \emph{jidōsha} — \textquotedblleft{}an automobile; a motor car\textquotedblright{}.

\begin{quote}
\textbf{\ja{ちちの} \ja{じどうしゃは} \ja{くろいです。}} — \emph{chichi no jidōsha wa kuroi desu} — \textquotedblleft{}My father's car is black.\textquotedblright{}
\end{quote}

\begin{cousinweb}
\textbf{\ja{じどうしゃ}} (\emph{jidōsha}) is Sino-Japanese (kango): three Chinese characters read with their on'yomi, readings borrowed from Chinese. They are \emph{ji}, \textquotedblleft{}self\textquotedblright{} (Middle Chinese \emph{dzijH}), \emph{dō}, \textquotedblleft{}to move\textquotedblright{} (\emph{duwngX}) and \emph{sha}, \textquotedblleft{}cart\textquotedblright{} (\emph{tsyhae}) (KANJIDIC2; Middle Chinese in Baxter's transcription).
\end{cousinweb}

\begin{grammarlens}[title={What you've built}]
One more word for everyday life at A2.
\end{grammarlens}

\subsection*{Your turn}
\begin{itemize}
\raggedright
  \item \emph{Say it:} \emph{jidōsha}
  \item \emph{Say it:} \emph{jidōsha}, once more
  \item \emph{Say it:} \emph{chichi no jidōsha wa kuroi desu} — My father's car is black.
  \item \emph{Recall:} say the Japanese for to rush out, then the Japanese for to bump into, then say \emph{jidōsha} again
\end{itemize}

\begin{tcolorbox}[breakable,colback=teal!4,colframe=teal!35!black,title={Before you move on}]
What does \ja{じどうしゃ} mean? (\textquotedblleft{}An automobile; a motor car\textquotedblright{}.) Say it once more.
\end{tcolorbox}
\section[\texorpdfstring{\ja{うんてん} (unten)}{うんてん (unten)}]{\ja{うんてん} (unten) — driving; operating (a machine)}

\label{lesson:JA-C202-unten}



\begin{quote}
Before the new one: say the Japanese for to bump into, then the Japanese for an automobile.
\end{quote}

\subsection*{You'll want to know: \ja{うんてん}}
\textbf{\ja{うんてん}} — \emph{unten} — \textquotedblleft{}driving; operating (a machine)\textquotedblright{}.

\begin{quote}
\textbf{\ja{うんてんは} \ja{たのしいです。}} — \emph{unten wa tanoshii desu} — \textquotedblleft{}Driving is fun.\textquotedblright{}
\end{quote}

\begin{cousinweb}
\textbf{\ja{うんてん}} (\emph{unten}) is Sino-Japanese (kango): two Chinese characters read with their on'yomi, readings borrowed from Chinese. They are \emph{un}, \textquotedblleft{}to carry\textquotedblright{} (Middle Chinese \emph{hjunH}) and \emph{ten}, \textquotedblleft{}to turn\textquotedblright{} (\emph{trjwenX}) (KANJIDIC2; Middle Chinese in Baxter's transcription).
\end{cousinweb}

\begin{grammarlens}[title={What you've built}]
One more word for everyday life at A2.
\end{grammarlens}

\subsection*{Your turn}
\begin{itemize}
\raggedright
  \item \emph{Say it:} \emph{unten}
  \item \emph{Say it:} \emph{unten}, once more
  \item \emph{Say it:} \emph{unten wa tanoshii desu} — Driving is fun.
  \item \emph{Recall:} say the Japanese for to bump into, then the Japanese for an automobile, then say \emph{unten} again
\end{itemize}

\begin{tcolorbox}[breakable,colback=teal!4,colframe=teal!35!black,title={Before you move on}]
What does \ja{うんてん} mean? (\textquotedblleft{}Driving; operating (a machine)\textquotedblright{}.) Say it once more.
\end{tcolorbox}
\section[\texorpdfstring{\ja{うんてんしゅ} (untenshu)}{うんてんしゅ (untenshu)}]{\ja{うんてんしゅ} (untenshu) — a driver}

\label{lesson:JA-C202-untenshu}



\begin{quote}
Before the new one: say the Japanese for an automobile, then the Japanese for driving.
\end{quote}

\subsection*{You'll want to know: \ja{うんてんしゅ}}
\textbf{\ja{うんてんしゅ}} — \emph{untenshu} — \textquotedblleft{}a driver\textquotedblright{}.

\begin{quote}
\textbf{\ja{ちちは} \ja{うんてんしゅです。}} — \emph{chichi wa untenshu desu} — \textquotedblleft{}My father is a driver.\textquotedblright{}
\end{quote}

\begin{cousinweb}
\textbf{\ja{うんてんしゅ}} (\emph{untenshu}) is Sino-Japanese (kango): three Chinese characters read with their on'yomi, readings borrowed from Chinese. They are \emph{un}, \textquotedblleft{}to carry\textquotedblright{} (Middle Chinese \emph{hjunH}), \emph{ten}, \textquotedblleft{}to turn\textquotedblright{} (\emph{trjwenX}) and \emph{shu}, \textquotedblleft{}hand\textquotedblright{} (\emph{syuwX}) (KANJIDIC2; Middle Chinese in Baxter's transcription).
\end{cousinweb}

\begin{grammarlens}[title={What you've built}]
One more word for everyday life at A2.
\end{grammarlens}

\subsection*{Your turn}
\begin{itemize}
\raggedright
  \item \emph{Say it:} \emph{untenshu}
  \item \emph{Say it:} \emph{untenshu}, once more
  \item \emph{Say it:} \emph{chichi wa untenshu desu} — My father is a driver.
  \item \emph{Recall:} say the Japanese for an automobile, then the Japanese for driving, then say \emph{untenshu} again
\end{itemize}

\begin{tcolorbox}[breakable,colback=teal!4,colframe=teal!35!black,title={Before you move on}]
What does \ja{うんてんしゅ} mean? (\textquotedblleft{}A driver\textquotedblright{}.) Say it once more.
\end{tcolorbox}
\section[\texorpdfstring{\ja{とめる} (tomeru)}{とめる (tomeru)}]{\ja{とめる} (tomeru) — to halt (a car); to park}

\label{lesson:JA-C202-tomeru}



\begin{quote}
Before the new one: say the Japanese for driving, then the Japanese for a driver.
\end{quote}

\subsection*{You'll want to know: \ja{とめる}}
\textbf{\ja{とめる}} — \emph{tomeru} — \textquotedblleft{}to halt (a car); to park\textquotedblright{}.

\begin{quote}
\textbf{\ja{くるまを} \ja{ここに} \ja{とめます。}} — \emph{kuruma o koko ni tomemasu} — \textquotedblleft{}I park the car here.\textquotedblright{}
\end{quote}

Polite \textbf{\ja{とめます}}, \emph{tomemasu}; an ichidan verb. Its partner is \textbf{\ja{とまる}}, \emph{tomaru}, \textquotedblleft{}to come to a stop\textquotedblright{}.

\begin{cousinweb}
\textbf{\ja{とめる}} (\emph{tomeru}) is a native Japanese word (wago), inherited rather than borrowed from Chinese; its kanji was attached later as a kun'yomi, and it pairs with \emph{tomaru}, \textquotedblleft{}to come to a stop\textquotedblright{} (JMdict, KANJIDIC2).
\end{cousinweb}

\begin{grammarlens}[title={What you've built}]
One more word for everyday life at A2.
\end{grammarlens}

\subsection*{Your turn}
\begin{itemize}
\raggedright
  \item \emph{Say it:} \emph{tomeru}
  \item \emph{Say it:} \emph{tomeru}, once more
  \item \emph{Say it:} \emph{kuruma o koko ni tomemasu} — I park the car here.
  \item \emph{Recall:} say the Japanese for driving, then the Japanese for a driver, then say \emph{tomeru} again
\end{itemize}

\begin{tcolorbox}[breakable,colback=teal!4,colframe=teal!35!black,title={Before you move on}]
What does \ja{とめる} mean? (\textquotedblleft{}To halt (a car); to park\textquotedblright{}.) Say it once more.
\end{tcolorbox}
\section[\texorpdfstring{\ja{ちゅうしゃじょう} (chūshajō)}{ちゅうしゃじょう (chūshajō)}]{\ja{ちゅうしゃじょう} (chūshajō) — a car park}

\label{lesson:JA-C202-chuushajou}



\begin{quote}
Before the new one: say the Japanese for a driver, then the Japanese for to halt.
\end{quote}

\subsection*{You'll want to know: \ja{ちゅうしゃじょう}}
\textbf{\ja{ちゅうしゃじょう}} — \emph{chūshajō} — \textquotedblleft{}a car park\textquotedblright{}.

\begin{quote}
\textbf{\ja{ちゅうしゃじょうは} \ja{えきの} \ja{となりです。}} — \emph{chūshajō wa eki no tonari desu} — \textquotedblleft{}The car park is next to the station.\textquotedblright{}
\end{quote}

\begin{cousinweb}
\textbf{\ja{ちゅうしゃじょう}} (\emph{chūshajō}) is Sino-Japanese (kango): three Chinese characters read with their on'yomi, readings borrowed from Chinese. They are \emph{chū}, \textquotedblleft{}to stop over\textquotedblright{} (Middle Chinese \emph{trjuH}), \emph{sha}, \textquotedblleft{}cart\textquotedblright{} (\emph{tsyhae}) and \emph{jō}, \textquotedblleft{}place\textquotedblright{} (\emph{drjang}) (KANJIDIC2; Middle Chinese in Baxter's transcription).
\end{cousinweb}

\begin{grammarlens}[title={What you've built}]
One more word for everyday life at A2.
\end{grammarlens}

\subsection*{Your turn}
\begin{itemize}
\raggedright
  \item \emph{Say it:} \emph{chūshajō}
  \item \emph{Say it:} \emph{chūshajō}, once more
  \item \emph{Say it:} \emph{chūshajō wa eki no tonari desu} — The car park is next to the station.
  \item \emph{Recall:} say the Japanese for a driver, then the Japanese for to halt, then say \emph{chūshajō} again
\end{itemize}

\begin{tcolorbox}[breakable,colback=teal!4,colframe=teal!35!black,title={Before you move on}]
What does \ja{ちゅうしゃじょう} mean? (\textquotedblleft{}A car park\textquotedblright{}.) Say it once more.
\end{tcolorbox}
